% New standalone review increment. No existing source is modified or merged.
\documentclass[11pt]{article}
\usepackage[T1]{fontenc}\usepackage[utf8]{inputenc}\usepackage{lmodern}
\usepackage[margin=.9in]{geometry}
\usepackage{amsmath,amssymb,amsthm,mathtools,microtype,xurl}
\usepackage[hidelinks]{hyperref}\usepackage{fancyhdr}
\setlength{\headheight}{15pt}\setlength{\emergencystretch}{3em}
\allowdisplaybreaks[2]\numberwithin{equation}{section}
\newtheorem{theorem}{Theorem}[section]
\newtheorem{proposition}[theorem]{Proposition}
\newtheorem{lemma}[theorem]{Lemma}\newtheorem{corollary}[theorem]{Corollary}
\theoremstyle{remark}\newtheorem{remark}[theorem]{Remark}
\newcommand{\E}{\mathbb E}\newcommand{\dd}{\,d}
\newcommand{\Qtwo}{\mathcal Q_2}\newcommand{\Clk}{\mathcal C_q}
\newcommand{\Sec}{\mathcal S_q}
\pagestyle{fancy}\fancyhf{}\fancyfoot[C]{\thepage}
\fancyhead[L]{\small AN03: bulk clock and Q2 energy}
\fancyhead[R]{\small Review increment v1}
\title{A separate bulk clock, the secondary Q2 tail,\\and regenerated transverse energy}
\author{Analytical increment for independent review}\date{Version 1}
\begin{document}\maketitle
\begin{abstract}
We separate the cohort used to control the exact response clock from the population whose output must be retained. An exact cohort-change identity explains the logarithmic defect generated by a large, slowly growing amplitude floor. A fixed target-only Q2 subcohort with initial $-\cos\alpha\ge q^{1/4}$ has the bulk seed scale $q^{10(1-\lambda)}$ and stays uniformly outside the strong noise-scale chart through the strong comparison clock. A Cartesian ratio bootstrap bounds the original-flow bulk clock defect from explicit residual norms and entry ratios. Under the same pre-exit distance and capped radial-gain inputs proposed for the right-half-circle complement, the recruited Q2 secondary tail has energy at most $C[q^5+q^{\Gamma_\nu}H]$, where $\Gamma_\nu=\gamma-\nu[10(1-\lambda)-2]>3/10$ for $\nu\le1/100$. Finally, exact Cartesian energy inequalities propagate a transverse budget $J\le C[q^3+q^2H]$, rather than the generally unsuitable $J\le Cq^3$, and recover the desired integrated forcing using the separate bulk clock. Target-only conclusions and conditional original-flow implications are distinguished throughout. The initialized radial-gain, profile-transfer, and residual-norm estimates are not proved here.
\end{abstract}

\section{Setting, source contracts, and energy variables}
Use the normalized laws
\[
 P_1=\mathcal N(e_1,q^2I_2),\quad P_2=\mathcal N(\rho e_2,q^2I_2),\quad
 P=(P_1+P_2)/2,\quad \lambda=\rho^2\in[169/400,49/100].
\]
The fixed initial label law is $\dd\lambda_0=\dd\alpha/(2\pi)$ and
$U(0,\alpha)=W(0,\alpha)=\sqrt{S_0}(\cos\alpha,\sin\alpha)$.
For the original residual $e(X)=X-f(X)$,
\[
 R_\theta=\E_P[e(X)X^T\mathbf1_{\{a_\theta^TX>0\}}],\qquad
 U'=R_\theta^TW,\quad W'=R_\theta U.
\]
Balance $|U|=|W|=\sqrt m$ holds, but alignment need not hold. For a fixed cohort $G$, set
\begin{align}
 H_G&=\frac14\int_G(U_2+W_2)^2\dd\lambda_0,&
 A_G&=\frac14\int_G(U_2-W_2)^2\dd\lambda_0,\notag\\
 E_G&=H_G+A_G=\frac12\int_G(U_2^2+W_2^2)\dd\lambda_0,&
 J_G&=\frac12\int_G(U_1^2+W_1^2)\dd\lambda_0.
 \label{inc:AN03:bulkclock:v1:energies}
\end{align}
The response and its cohort-independent clock are
\[
 Q=\lambda+q^2,\qquad c(t)=\E_{P_2}[X_2f_2(X)]/Q,\qquad
 I(t)=Q\int_{t_0}^t(1-c(v))\dd v.
\]
Source CLOCK proves, for each fixed cohort with positive $H_G$,
\begin{equation}\label{inc:AN03:bulkclock:v1:clock}
 H_G'=I'H_G+\mathcal D_G,\qquad
 \varepsilon_G=\mathcal D_G/H_G,\qquad
 \int_s^t\varepsilon_G=\log\frac{H_G(t)}{H_G(s)}-[I(t)-I(s)].
\end{equation}
It also proves the exact original-state bound
\begin{equation}\label{inc:AN03:bulkclock:v1:cross}
 \mathfrak F_G:=\frac{\|(f_G)_2\|_{L^2(P_1)}+\|(f_G)_1\|_{L^2(P_2)}}q
 \le8\left[E_G+J_G+\frac{\sqrt{E_GJ_G}}q\right].
\end{equation}
All time-dependent groups below are actually fixed sets of initial labels, possibly selected using an analytically specified comparison at one deterministic time. No moving-mask derivative is used.

\section{Changing the clock cohort: an exact identity and the amplitude floor}
\begin{proposition}[Cohort change]\label{inc:AN03:bulkclock:v1:cohortchange}
Let $C$ and $S$ be disjoint fixed cohorts, $H_C>0$, and $H_{C\cup S}=H_C+H_S$. Then
\begin{equation}\label{inc:AN03:bulkclock:v1:ratioidentity}
 \int_s^t(\varepsilon_{C\cup S}-\varepsilon_C)
 =\log\frac{1+H_S(t)/H_C(t)}{1+H_S(s)/H_C(s)}.
\end{equation}
In particular, for every $r>0$,
\[
 \left(\frac{H_{C\cup S}(t)}{H_{C\cup S}(s)}\right)^r
 e^{-r\int_s^t\varepsilon_{C\cup S}}
 =e^{r[I(t)-I(s)]}
 =\left(\frac{H_C(t)}{H_C(s)}\right)^r e^{-r\int_s^t\varepsilon_C}.
\]
\end{proposition}
\begin{proof}
Subtract \eqref{inc:AN03:bulkclock:v1:clock} for the two positive amplitudes. The last display follows by exponentiation. No positivity of individual label amplitudes and no monotonicity of $H_G$ are required.
\end{proof}

For a precise scalar illustration, let $I(t_0)=0$, $I'\ge0$, and
\[
 H_C=b e^{I},\qquad H_S=a,\qquad a,b>0.
\]
Then $\varepsilon_C=0$ and
\begin{equation}\label{inc:AN03:bulkclock:v1:floormodel}
 \varepsilon_{C\cup S}=-I'\frac{a}{a+b e^I},\qquad
 \int_{t_0}^t\varepsilon_{C\cup S}
 =\log\frac{a e^{-I(t)}+b}{a+b}.
\end{equation}
If $a\asymp q^5$, $b\asymp q^k$, $k>5$, an interval reaching $I=(k-5)\log(1/q)+O(1)$ has defect $-(k-5)\log(1/q)+O(1)$. At $k=10(1-\lambda)$,
\[
 k-5=5-10\lambda\in[1/10,31/40].
\]
Thus an endpoint passage multiplier can retain the correct bulk power, whereas a forcing estimate which simply multiplies a fixed terminal amplitude by $e^{D_{\rm osc}}$ loses a negative power of $q$. This is an exact calculation for the stated scalar functions, not a proof that the initialized residual flow has a frozen $H_S$. More generally, \eqref{inc:AN03:bulkclock:v1:ratioidentity} proves the same relative defect loss whenever $H_S/H_C$ decreases from order $q^{5-k}$ to order one and the bulk defect is bounded on that interval.

\section{A fixed target-only clock subcohort and the intermediate band}
In this section only, hats denote the aligned target-only comparison
$\widehat U'=\E[X(\widehat U^TX)_+]$, $\widehat W=\widehat U$.
Write $\Qtwo=[\pi/2,\pi]$, $c_0=-\cos\alpha$, $s_0=\sin\alpha$,
$T=10\log(1/q)$, and $S_0=q^{10}$. Define
\[
 \gamma(\lambda)=1-5\lambda+\frac{2\lambda}{1-\lambda},\qquad
 \Clk=\{\alpha\in\Qtwo:c_0\ge q^{1/4}\}.
\]
Source Q2E proves the uniform crossing formula and crossing amplitude
\begin{equation}\label{inc:AN03:bulkclock:v1:crossinput}
 T_\times=\frac2\lambda\log\frac{c_0+q}{q(s_0+q)}+O(1),\qquad
 \widehat U_2(T_\times)\asymp\sqrt{S_0}\frac{c_0+q}{q}.
\end{equation}
It also proves the post-crossing bound
\begin{equation}\label{inc:AN03:bulkclock:v1:U1input}
 0\le\widehat U_1(T_\times+v)
 \le Cq\widehat U_2(T_\times)e^{v/2}\qquad(0\le v\le T-T_\times).
\end{equation}
These are uniform even when the original label approaches either Q2 endpoint.

\begin{theorem}[A bulk seed without the captured strip]\label{inc:AN03:bulkclock:v1:bulkseed}
Every label in $\Clk$ has crossed before $T$. For $T_\times\le t\le T$,
\begin{equation}\label{inc:AN03:bulkclock:v1:gap}
 q\cot\widehat\theta(t,\alpha)\le Cq^{\delta(\lambda)},\qquad
 \delta(\lambda)=\frac{1-\lambda}{\lambda}\left(\gamma(\lambda)-\frac14\right)
 =5\lambda-\frac{15}{4}+\frac{3}{4\lambda}.
\end{equation}
Here
\[
 \frac18<\frac{1861}{13520}\le\delta(\lambda)\le\frac{113}{490}<\frac14.
\]
Consequently $\sin\widehat\theta/q\ge c q^{-1/8}$ on this post-crossing interval; in particular these labels are not in a fixed strong noise-scale chart at $T$. Uniformly on $\Clk$,
\begin{equation}\label{inc:AN03:bulkclock:v1:pointseed}
 \widehat U_2(T,\alpha)^2\asymp S_0e^{\lambda T}(s_0+q)^2,
 \qquad \widehat H_{\Clk}(T)\asymp q^{10(1-\lambda)}.
\end{equation}
There is a non-clock, non-captured band
\[
 \mathcal M_q=\{\alpha\in\Qtwo:\tfrac12 q^{1/4}\le c_0<q^{1/4}\}
\]
with
\begin{equation}\label{inc:AN03:bulkclock:v1:middleband}
 \widehat H_{\mathcal M_q}(T)\asymp q^{10(1-\lambda)+1/4}.
\end{equation}
\end{theorem}
\begin{proof}
All target-only Q2 labels cross by $T$, because source Q2E gives maximum crossing coefficient $4/\lambda\le1600/169<10$. After crossing, the exact coordinate identity gives
\[
 \frac{\widehat U_2'}{\widehat U_2}
 =\frac\lambda2+d_q+\frac{\rho q}{2\sin\widehat\theta}
 K(-\rho\sin\widehat\theta/q),\qquad 0\le d_q\le q^2.
\]
The extra terms are nonnegative. Dividing \eqref{inc:AN03:bulkclock:v1:U1input} by the resulting lower exponential for $\widehat U_2$ gives
\[
 q\cot\widehat\theta(t)\le C q^2e^{\omega(t-T_\times)},\qquad
 \omega=(1-\lambda)/2.
\]
Insert \eqref{inc:AN03:bulkclock:v1:crossinput}, $t\le T$, $c_0\ge q^{1/4}$ and $s_0+q\le2$. The resulting exponent is exactly $\delta(\lambda)$ in \eqref{inc:AN03:bulkclock:v1:gap}. Its derivative is $5-3/(4\lambda^2)>0$ on the interval. Direct rational endpoint evaluation gives the displayed bounds. Since
\[
 \frac{\sin\widehat\theta}{q}=
 \frac1{\sqrt{q^2+(q\cot\widehat\theta)^2}},
\]
we get the claimed projected-margin bound. In particular, the integral of the extra Gaussian rate above is at most
\[
 C T q^{1/8}K(-c q^{-1/8})=o(1),
\]
and $\int d_q\le q^2T=o(1)$. Here $K(-x)\le\phi(x)$ for $x\ge0$ suffices. Thus
\[
 \widehat U_2(T)=\widehat U_2(T_\times)
 e^{\lambda(T-T_\times)/2}(1+o(1)).
\]
Using \eqref{inc:AN03:bulkclock:v1:crossinput} proves the pointwise comparison. Its integral over $\Clk$ is comparable to $S_0e^{\lambda T}$: the integral of $(s_0+q)^2$ is bounded above uniformly and below on a fixed interior Q2 subinterval contained in $\Clk$.

All the same estimates hold with $c_0\ge\tfrac12q^{1/4}$, changing only constants. On $\mathcal M_q$, $s_0\asymp1$ and the label measure is comparable to $q^{1/4}$. This proves \eqref{inc:AN03:bulkclock:v1:middleband}; the angular margin shows this band is not strong-chart captured. Alternatively source Q2E bounds all captured labels by $c_0\le Cq^\gamma$, and $\gamma>1/4$ uniformly.
\end{proof}

\begin{remark}[Two meanings of bulk must not be identified]
\label{inc:AN03:bulkclock:v1:twobulks}
The clock subcohort need not equal the whole primary weak-energy cohort. If $E_{\rm bulk}$ means the energy of $\Clk$ alone, the secondary-tail calculation below does not bound every omitted label by $O(q^\gamma E_{\rm bulk})$. The intermediate band in \eqref{inc:AN03:bulkclock:v1:middleband} has relative energy of order $q^{1/4}$ at entry. At that entry clock it is smaller than the separate $q^5$ floor, so this observation does not itself contradict a bound containing that floor. But its subsequent primary growth must be included: a transported $O(q^{1/4}H)$ term, or a larger definition of primary bulk with its own comparison to the clock cohort, is needed. Target-only entry scaling does not prove its original-flow transport. Corollary \ref{inc:AN03:bulkclock:v1:bulkweights} below supplies conditional transport when its original-field inputs hold on the relevant labels.
\end{remark}

\section{A bulk-clock defect bound from original residual norms}
The clock cohort can be controlled without assuming its signed defect in advance. The following conditional original-flow theorem uses Cartesian entry ratios and residual norms; the initialized verification of those hypotheses remains separate.

\begin{theorem}[Bulk ratio bootstrap and subinterval clock control]
\label{inc:AN03:bulkclock:v1:bulkclockclosure}
Let $C$ be a fixed cohort of the original balanced flow on $[t_0,t_b]$, with $0<H_C(t_0)<\infty$. Fix $0<\eta<1$, $c_b<1$, and constants independent of $q$. Assume
\begin{align}
 0\le c(t)&\le c_b,\qquad t_b-t_0\le L_T\log(1/q),\notag\\
 \|e_1\|_{L^2(P_1)}+\|e_2\|_{L^2(P_2)}&\le M,\qquad
 \|e_2\|_{L^2(P_1)}+\|e_1\|_{L^2(P_2)}\le Mq,
 \label{inc:AN03:bulkclock:v1:bulkresiduals}\\
 \int_{t_0}^{t_b}r(t)\dd t&\le B,\qquad
 r(t)=\operatorname*{ess\,sup}_{\alpha\in C}|(R_{\theta_\alpha})_{11}|.
 \notag
\end{align}
For a label set
\[
 z=\frac{U_2+W_2}{2},\qquad a=\frac{U_2-W_2}{2},\qquad
 j=\sqrt{\frac{U_1^2+W_1^2}{2}}.
\]
Suppose at entry, uniformly over $C$ up to null sets,
\begin{equation}\label{inc:AN03:bulkclock:v1:bulkentryratios}
 z(t_0)>0,\qquad \frac{|a(t_0)|}{z(t_0)}\le\frac14,\qquad
 \frac{qj(t_0)}{z(t_0)}\le C_0q^\eta.
\end{equation}
For sufficiently small $q$, $z>0$ throughout the interval, both $U_2,W_2$ stay positive, and, writing $\tau=t-t_0$,
\begin{align}
 \frac{qj(t)}{z(t)}&\le C[q^\eta e^{-b_*\tau/2}+q^2],\notag\\
 \frac{a_{\theta,2}(t)}q&\ge c_*q^{-\eta},\qquad
 b_*:=\frac{Q(1-c_b)}2.
 \label{inc:AN03:bulkclock:v1:bulkratioconclusion}
\end{align}
The exact individual and aggregate amplitude defects satisfy
\begin{equation}\label{inc:AN03:bulkclock:v1:bulkdefectconclusion}
 \int_s^t|\varepsilon_C(v)|\dd v
 \le C\left[q^\eta e^{-b_*(s-t_0)/2}+q^2(t-s)\right]
 \qquad(t_0\le s\le t\le t_b).
\end{equation}
The same bound holds for each individual defect
$\varepsilon_\alpha=(\log z_\alpha^2)'-Q(1-c)$.
In particular $D_{\rm osc}\le C[q^\eta+q^2\log(1/q)]=o(1)$ for this cohort. The target-only cutoff theorem supplies \eqref{inc:AN03:bulkclock:v1:bulkentryratios} with $a=0$ and $\eta=1/8$ for its comparison state; it does not itself transfer those Cartesian ratios to the original entry state.
\end{theorem}
\begin{proof}
Let $b(t)=Q(1-c(t))/2\ge b_*$. Cauchy--Schwarz in the two cluster expectations, as detailed later in Proposition \ref{inc:AN03:bulkclock:v1:residualinputs}, gives $|R_{12}|,|R_{21}|\le Lq$, with fixed $L\ge1$.
For any receiver whose second input component is positive,
\[
 R_{22}-b=\frac12\E_{P_1}[e_2X_2\mathbf1_{\{a_\theta^TX>0\}}]
 -\frac12\E_{P_2}[e_2X_2\mathbf1_{\{a_\theta^TX\le0\}}].
\]
The first term is $O(q^2)$ by \eqref{inc:AN03:bulkclock:v1:bulkresiduals}. For the other term, two Cauchy--Schwarz bounds give
\[
 |\E_{P_2}[e_2X_2\mathbf1_{\{a_\theta^TX\le0\}}]|
 \le \|e_2\|_{L^2(P_2)}(\E_{P_2}X_2^4)^{1/4}
       \mathbb P_{P_2}(a_\theta^TX\le0)^{1/4}.
\]
The fourth moment is uniformly bounded. The one-dimensional Gaussian projection has mean $\rho a_{\theta,2}$ and variance $q^2$, and its lower-tail probability is at most $\exp[-\rho^2a_{\theta,2}^2/(2q^2)]$. Hence
\begin{equation}\label{inc:AN03:bulkclock:v1:R22tail}
 |R_{22}-b|\le Cq^2+C\exp[-\rho^2a_{\theta,2}^2/(8q^2)].
\end{equation}

While $z>0$, put $x=qj/z$, $y=|a|/z$. The exact coordinate equations imply, with $\Delta=|R_{22}-b|$,
\begin{align*}
 z'/z&\ge b-\Delta-Lx,\\
 D^+x&\le(r-b+\Delta+Lx)x+Lq^2\sqrt{1+y^2},\\
 D^+y&\le(-2b+2\Delta+Lx)y+Lx.
\end{align*}
The norm derivative at $j=0$ is interpreted without dividing by $j$; the quotient divides only by the guarded positive $z$.

Choose
\[
 X_q=2e^BC_0q^\eta+\frac{8Le^B}{b_*}q^2
\]
and use the first-exit guards $x<X_q$, $y<1/2$, $z>0$. For small $q$, $X_q<b_*/(4L)$. Under these guards, $U_2=z+a\ge z/2$ and
\[
 |q\cot\theta|=q|U_1|/U_2\le2\sqrt2x,\qquad
 \frac{a_{\theta,2}}q=\frac1{\sqrt{q^2+(q\cot\theta)^2}}
 \ge c q^{-\eta}.
\]
Therefore \eqref{inc:AN03:bulkclock:v1:R22tail} gives $\Delta\le Cq^2$ throughout the guarded segment; in particular $\Delta\le b_*/8$ for small $q$. This is a consequence of the guards, not a separate assumed gate estimate. The inequalities reduce to
\[
 D^+x\le(r-b_*/2)x+2Lq^2,\qquad
 D^+y\le-3b_*y/2+Lx,\qquad z'/z\ge b_*/2.
\]
Variation of constants and $\int r\le B$ give
\[
 x(t)\le e^BC_0q^\eta e^{-b_*\tau/2}
                +\frac{4Le^B}{b_*}q^2\le X_q/2.
\]
The $y$ bound is consequently
\[
 y(t)\le\frac14+\frac{LX_q}{3b_*}\le\frac13<\frac12.
\]
The positive lower logarithmic rate of $z$ prevents a zero. All guards close with slack, establishing the ratio and gate conclusions.

Finally, the exact $z$ equation gives
\[
 |(\log z^2)'-2b|\le2\Delta+2Lx
 \le C[q^\eta e^{-b_*\tau/2}+q^2].
\]
Integrating yields the individual defect bound. The aggregate defect is the current $z_\alpha^2$-weighted average of the individual defects, so the same uniform pointwise bound and its subinterval integral apply to $\varepsilon_C$. No fixed normalized radial weights are assumed.
\end{proof}

\begin{corollary}[Relative amplitude of bulk subsets]
\label{inc:AN03:bulkclock:v1:bulkweights}
Suppose Theorem \ref{inc:AN03:bulkclock:v1:bulkclockclosure} holds on a fixed group containing positive-amplitude subsets $G_1,G_2$. Put $D_q=C[q^\eta+q^2\log(1/q)]$. Then
\[
 e^{-2D_q}\le
 \frac{H_{G_2}(t)/H_{G_1}(t)}{H_{G_2}(t_0)/H_{G_1}(t_0)}
 \le e^{2D_q}.
\]
If the original entry amplitudes of the band $\mathcal M_q$ and $\Clk$ have the ratio $\asymp q^{1/4}$, and the bootstrap inputs hold on their union, that amplitude ratio stays $\asymp q^{1/4}$ throughout the interval. Their raw energies satisfy $H_G\le E_G\le(10/9)H_G$ under the same bootstrap. This is a conditional original-flow transport statement, not a transfer of the target-only entry profile.
\end{corollary}
\begin{proof}
The uniform individual defect bound also bounds every subset's aggregate defect. Subtract the two logarithmic clock identities and exponentiate. The energy comparison follows from the proof's $|a|/z\le1/3$.
\end{proof}

\begin{remark}
This theorem closes the bulk clock and its input gate margin from explicit original-field hypotheses. It does not prove those residual norms or the integrated $R_{11}$ estimate from isotropic initialization. Nor does the displayed projected margin alone give bounded weak noise-scaled coordinates at $t_b$. A relative Cartesian entry comparison is needed: an unspecified $o(1)$ angular error does not automatically preserve \eqref{inc:AN03:bulkclock:v1:bulkentryratios} for labels with small second-coordinate amplitude.
\end{remark}

\section{Secondary-tail moments, capped gain, and the exact exponent}
Let $a_\rho=\rho K(\rho a_\rho)$ and write
\[
 P_\rho=\Phi(\rho a_\rho),\quad d=(1-\lambda P_\rho)/2,\quad
 s=\frac{2d}{1-\lambda},\quad
 h_1=q^{\kappa d-2s},\quad p_2=\frac{1-3\lambda/2}{d}.
\]
At $\kappa=10$, source Q2E supplies the finite-cutoff recruited tail, with $0<p_2<1$. Its initial radial submeasure, after a bounded strong-mass normalization, obeys
\begin{equation}\label{inc:AN03:bulkclock:v1:tailcontract}
 \eta([0,Z])\le Cq^3,\quad \operatorname{supp}\eta\subset[0,Z],\quad
 \eta\{z>u\}\le Cq^3\min\{1,(h_1/u)^{p_2}\}.
\end{equation}
The fixed $Z$ is a chart-entry displacement cap, not a global untruncated power-law hypothesis. The upper bound follows from Q2E's full weight/displacement profiles; the two-sided tail there was asserted only away from the final cutoff. Transfer of this contract to the original joint input/output displacement remains an assumption.

\begin{lemma}[High moments of the truncated secondary tail]
\label{inc:AN03:bulkclock:v1:moments}
Under \eqref{inc:AN03:bulkclock:v1:tailcontract}, for $r>p_2$,
\[
 \int z^r\dd\eta\le C\left(1+\frac{r}{r-p_2}\right)
 q^3h_1^{p_2}Z^{r-p_2}.
\]
In particular this holds uniformly for $r=2$ and $r=2(1+\nu)$, $0\le\nu\le1/100$. If the source's two-sided tail holds at a fixed positive displacement below $Z$, these two moments are also bounded below by a constant times $q^3h_1^{p_2}$.
\end{lemma}
\begin{proof}
Tonelli gives $\int z^r\dd\eta=r\int_0^Z u^{r-1}\eta\{z>u\}\dd u$.
If $h_1\ge Z$, the conclusion follows directly from the mass bound. Otherwise split at $h_1$ and integrate the displayed powers. For the lower bound, use a fixed threshold $u_*>0$ with $\int z^r\dd\eta\ge u_*^r\eta\{z>u_*\}$. For the stated moment orders $r-p_2>1$, so all upper constants are uniform on the compact parameter range.
\end{proof}

The next proposition makes explicit the angular input needed to turn radial mass into weak-coordinate energy. It uses the same distance envelope involved in deriving the proposed capped radial-gain law; it does not infer energy from mass alone.

\begin{proposition}[From pre-exit distance and capped mass gain to energy]
\label{inc:AN03:bulkclock:v1:gainenergy}
Let $\Sec$ be fixed labels whose entry radial measure is $\dd\eta=m(t_0,\alpha)\dd\lambda_0$ and whose entry joint displacement is $z_\alpha\le Z$. Let the same-field reference have both strong chart coordinates bounded by $R$. Define $D_\alpha$ to be the sum of absolute joint coordinate differences while $qD_\alpha\le\delta_0$. Suppose $I(t_0)=0$, $I'\ge0$, $0<\beta\le1$, and, until the first such chart exit,
\[
 D_\alpha(t)\le C_Dz_\alpha e^{\beta I(t)/2}.
\]
Suppose the actual mass satisfies at every time considered
\begin{equation}\label{inc:AN03:bulkclock:v1:gaincontract}
 m(t,\alpha)\le C_m m(t_0,\alpha)
 \max\{1,C_gq^2z_\alpha^2e^{\beta I(t)}\}^{1+\nu},\qquad \nu\ge0.
\end{equation}
Then, with $\mathcal A(t)=q^2e^{\beta I(t)}$ and constants allowed to depend on fixed $\delta_0,R,C_D,C_m,C_g,Z$,
\begin{equation}\label{inc:AN03:bulkclock:v1:labelenergy}
 E_{\Sec}(t)\le Cq^2\eta([0,Z])+
 C\mathcal A(t)\int z^2\dd\eta+
 C\mathcal A(t)^{1+\nu}\int z^{2(1+\nu)}\dd\eta.
\end{equation}
No bound on post-exit chart coordinates and no assertion of completed weak return are needed for this implication.
\end{proposition}
\begin{proof}
Fix a label and put $X=q^2z_\alpha^2e^{\beta I(t)}$. Choose a fixed $\eta_*>0$ small enough that $C_D\sqrt{\eta_*}<\delta_0/2$. At entry $qZ$ is small. As long as $X<\eta_*$, a first exit of the physical distance guard would contradict $qD_\alpha\le C_D\sqrt X<\delta_0/2$. Since $I$ is nondecreasing, the comparison threshold cannot have been crossed earlier and then undone. On this segment \eqref{inc:AN03:bulkclock:v1:gaincontract} bounds mass by a fixed multiple of its entry value. Also
\[
 \frac{U_2^2+W_2^2}{2}
 =\frac m2[\sin^2(qx)+\sin^2(qy)]
 \le C_Rm q^2(1+D_\alpha^2)
 \le C m(t_0,\alpha)(q^2+X).
\]
When $X\ge\eta_*$, whether or not an actual exit occurred, balance gives $(U_2^2+W_2^2)/2\le m$, and the mass bound is at most $C_{\eta_*}m(t_0,\alpha)X^{1+\nu}$. Combining the two alternatives and integrating proves \eqref{inc:AN03:bulkclock:v1:labelenergy}. No division by $z_\alpha$ is used; a zero-displacement label stays in the first alternative.
\end{proof}

\begin{corollary}[Secondary-tail amplitude margin]
\label{inc:AN03:bulkclock:v1:secondary}
Assume the preceding two results' contracts for the original-flow secondary cohort. Suppose a possibly different fixed clock cohort has amplitude $H>0$, entry value $H_0\asymp q^k$ with $k=\kappa(1-\lambda)$, and
\[
 \int_s^t\varepsilon_H\ge-D\quad(t_0\le s\le t\le t_b),\qquad
 H(t_b)\le H_b,
\]
where $D,H_b$ are independent of $q$, and $0\le c\le c_b<1$. Then
\begin{equation}\label{inc:AN03:bulkclock:v1:secondarydetail}
 E_{\Sec}(t)\le C\big[q^5+q^{\gamma_\kappa}H(t)
          +q^{\Gamma_\nu}H(t)^{1+\nu}\big],
\end{equation}
where the algebraic expressions are
\begin{align}
 \gamma_\kappa&=1-\frac\lambda2\left(\kappa-\frac4{1-\lambda}\right),\notag\\
 \Gamma_\nu&=\gamma_\kappa-\nu[\kappa(1-\lambda)-2].
 \label{inc:AN03:bulkclock:v1:gamma}
\end{align}
For $\kappa=10$, $0\le\nu\le1/100$,
\begin{equation}\label{inc:AN03:bulkclock:v1:gammamargin}
 \Gamma_\nu\ge\frac{289169}{924000}>\frac3{10},\qquad
 E_{\Sec}(t)\le C[q^5+q^{\Gamma_\nu}H(t)].
\end{equation}
The general-$\kappa$ expression is exponent algebra under the displayed profile contract, not a new target-only profile theorem at every $\kappa$.
\end{corollary}
\begin{proof}
The exact clock and its lower defect bound give $e^{I(t)}\le e^D H(t)/H_0$. Since $I\ge0$ and $\beta\le1$, $e^{\beta I}\le e^I$. Insert these inequalities and the two moment estimates into \eqref{inc:AN03:bulkclock:v1:labelenergy}. The linear-gain exponent is
\[
 3+2-k+p_2(\kappa d-2s)=\gamma_\kappa,
\]
and the power-gain exponent is this number minus $\nu(k-2)$. To verify the identity directly, use
\[
 \frac{\kappa d-2s}{b}=\frac\lambda2\left(\kappa-\frac4{1-\lambda}\right),
 \quad p_2b=\frac2\lambda-3,\quad b=\frac{2d}{\lambda}.
\]
Also the all-subinterval bound gives $H(t)\le e^D H_b$, so the last energy term is at most a fixed constant times $q^{\Gamma_\nu}H(t)$.

At $\kappa=10$, the derivative of $\gamma$ is $-5+2/(1-\lambda)^2>0$ on the entire ratio interval. Therefore
\[
 \gamma\ge\gamma(169/400)=\frac{6481}{18480}=\frac{32405}{92400}>\frac7{20}.
\]
Furthermore $k-2\le151/40$, whence
\[
 \Gamma_\nu\ge\frac{6481}{18480}-\frac{151}{4000}
 =\frac{289169}{924000}>\frac3{10}.
\]
Here $k-2>0$, so $q^{\gamma}H$ is absorbed by $q^{\Gamma_\nu}H$.
\end{proof}

\begin{remark}[The meaning of the sharp exponent]
For a genuinely quadratic gain ($\nu=0$), the secondary exponent is exactly $\gamma_\kappa$. The profile-variable identity $2b-a_2=2(1-P_\rho)>0$, $a_2=2/\lambda-2$, explains why the high moments are dominated by last arrivals. A fixed positive gain-power loss changes that exponent to $\Gamma_\nu$; it cannot be hidden in a constant. If the only loss is $\nu=O(q^2)$, its factor $q^{-\nu(k-2)}=\exp[O(q^2\log(1/q))]$ is uniformly bounded and tends to one. The identity $\gamma_\kappa>0\iff\kappa<2/\lambda+4/(1-\lambda)$ is algebraic. Crossing the latter target-only timing boundary does not by itself prove capture in the trained residual flow.
\end{remark}

\section{Regenerated transverse energy and a forcing theorem with a separate clock}
The raw bound $J_{\Qtwo}\lesssim q^3$ is appropriate at the target-only strong entry clock, but not as a general terminal weak-chart requirement. At a bounded weak-chart state,
\[
 U=\sqrt m(\sin(qu),\cos(qu)),\qquad
 W=\sqrt m(-\sin(qv),\cos(qv)),
\]
so
\begin{equation}\label{inc:AN03:bulkclock:v1:weakJ}
 J_G=\frac{q^2}{2}\int_Gm(u^2+v^2)\dd\lambda_0
       +O\left(q^4\int_Gm(u^4+v^4)\dd\lambda_0\right),\quad
 E_G=N_G+O(q^2N_G).
\end{equation}
For example a fixed positive fraction of weak mass with $v$ bounded away from zero has $J_G\gtrsim q^2N_G$. The selected weak branch approaches $v=a_\rho/\lambda>0$ at its small-mass resident end. Thus a fixed nonzero weak amplitude can regenerate $q^2H$ transverse energy; the correct useful allowance includes this term. This is a static chart calculation, not an original-flow capture assertion.

\begin{lemma}[Exact Cartesian transverse-energy inequality]
\label{inc:AN03:bulkclock:v1:transverse}
For any fixed original-flow cohort $G$, define
\[
 r_G(t)=\operatorname*{ess\,sup}_{\alpha\in G}|(R_{\theta_\alpha})_{11}|,\qquad
 \ell_G(t)=\operatorname*{ess\,sup}_{\alpha\in G}
       \max\{|(R_{\theta_\alpha})_{12}|,|(R_{\theta_\alpha})_{21}|\}.
\]
Then
\begin{equation}\label{inc:AN03:bulkclock:v1:jineq}
 D^+\sqrt{J_G}\le r_G\sqrt{J_G}+\ell_G\sqrt{E_G}.
\end{equation}
In particular, if $\int_{t_0}^{t_b}r_G\le B$,
\[
 \sqrt{J_G(t)}\le e^B\left[\sqrt{J_G(t_0)}+
       \int_{t_0}^t\ell_G(s)\sqrt{E_G(s)}\dd s\right].
\]
\end{lemma}
\begin{proof}
For a label, $U_1'=R_{11}W_1+R_{21}W_2$ and $W_1'=R_{11}U_1+R_{12}U_2$. Therefore
\[
 J_G'=\int_G[2R_{11}U_1W_1+R_{21}U_1W_2+R_{12}W_1U_2]\dd\lambda_0
 \le2r_GJ_G+2\ell_G\sqrt{J_GE_G}.
\]
Cauchy--Schwarz supplies the last inequality. The square-root estimate follows wherever $J_G>0$. At zero, use the upper norm derivative in the Hilbert space of $(U_1,W_1)/\sqrt2$, or regularize the norm and pass to the limit; no division by zero is needed. Variation of constants proves the conclusion. The derivatives and integrals are valid almost everywhere in the characteristic class with finite coordinate energies on compact intervals.
\end{proof}

\begin{theorem}[Separate-clock energy and forcing control]
\label{inc:AN03:bulkclock:v1:forcing}
Suppose $H$ is the positive amplitude of a fixed clock cohort, with
\[
 0\le c\le c_b<1,\quad b_0=Q(1-c_b)>0,\quad
 \int_s^t\varepsilon_H\ge-D\quad(t_0\le s\le t\le t_b).
\]
Let $G$ be any other fixed cohort; it need not equal or be contained in the clock cohort. Assume, uniformly on this interval,
\begin{equation}\label{inc:AN03:bulkclock:v1:forcinginputs}
 E_G(t)\le C_E[q^5+q^\sigma H(t)],\quad
 J_G(t_0)\le C_Jq^3,\quad
 \ell_G(t)\le Lq,\quad \int_{t_0}^{t_b}r_G\le B,
\end{equation}
where $\sigma\ge0$, all constants are independent of $q$, and
$t_b-t_0\le L_T\log(1/q)$. Then
\begin{align}
 J_G(t)&\le C[q^3+q^{2+\sigma}H(t)],
 \label{inc:AN03:bulkclock:v1:Jprop}\\
 \mathfrak F_G(t)&\le C[q^3+q^\sigma H(t)
                  +q^{(1+\sigma)/2}\sqrt{H(t)}].
 \label{inc:AN03:bulkclock:v1:Fpoint}
\end{align}
If $H(t_b)\le H_b$, then
\begin{equation}\label{inc:AN03:bulkclock:v1:Fintegral}
 \int_{t_0}^{t_b}\mathfrak F_G(t)\dd t
 \le C\left[q^3(t_b-t_0)+
 \frac{e^Dq^\sigma H_b}{b_0}+
 \frac{2e^{D/2}q^{(1+\sigma)/2}\sqrt{H_b}}{b_0}\right].
\end{equation}
The cases $\sigma=0$ and $\sigma=\Gamma_\nu$ give respectively the whole-Q2 and secondary-Q2 budgets, provided their other displayed inputs have been established.
\end{theorem}
\begin{proof}
The exact clock gives, for every $s\le t$,
\[
 H(s)\le e^DH(t)e^{-b_0(t-s)}.
\]
Thus $\int_{t_0}^t\sqrt H\le2e^{D/2}\sqrt{H(t)}/b_0$. Insert the energy assumption and $\ell_G\le Lq$ into Lemma \ref{inc:AN03:bulkclock:v1:transverse} to obtain
\[
 \sqrt{J_G(t)}\le C\left[q^{3/2}+q^{7/2}(t-t_0)
                  +q^{1+\sigma/2}\sqrt{H(t)}\right].
\]
The middle term is bounded by a constant times $q^{3/2}$ on the prescribed logarithmic interval. Squaring proves \eqref{inc:AN03:bulkclock:v1:Jprop}. Combining it with the energy input in the exact cross-output estimate \eqref{inc:AN03:bulkclock:v1:cross} gives \eqref{inc:AN03:bulkclock:v1:Fpoint}: expand the four products inside the square root, retaining the larger terms. In particular $\sqrt{q^\sigma H\,q^{2+\sigma}H}/q=q^\sigma H$, so regenerated transverse energy costs a term already present, not a new divergent power. Finally integrate the same backward bound on $H$ and its square root with terminal time $t_b$.
\end{proof}

\begin{proposition}[Concrete residual-norm inputs for the transverse theorem]
\label{inc:AN03:bulkclock:v1:residualinputs}
All receiver selectors satisfy
\begin{align*}
 |R_{12}|&\le\tfrac12[q\|e_1\|_{L^2(P_1)}+\sqrt Q\|e_1\|_{L^2(P_2)}],\\
 |R_{21}|&\le\tfrac12[\sqrt{1+q^2}\|e_2\|_{L^2(P_1)}+q\|e_2\|_{L^2(P_2)}],\\
 |R_{11}|&\le\tfrac12[\sqrt{1+q^2}\|e_1\|_{L^2(P_1)}+q\|e_1\|_{L^2(P_2)}].
\end{align*}
Consequently bounded diagonal residual norms and $O(q)$ cross residual norms imply $\ell_G\le Lq$ for every cohort. A sufficient integrated $r_G$ input is
\[
 \|e_1\|_{L^2(P_1)}\le a_s(t)+Cq^2,\quad
 \int_{t_0}^{t_b}a_s\le B_s,\quad
 \|e_1\|_{L^2(P_2)}\le Cq.
\]
These are original-flow hypotheses, not claims established by the target-only energy theorem.
\end{proposition}
\begin{proof}
Apply Cauchy--Schwarz to each of the two cluster terms in each entry of $R$, dropping its indicator inside the norm. Use $\|X_2\|_{L^2(P_1)}=\|X_1\|_{L^2(P_2)}=q$, $\|X_1\|_{L^2(P_1)}=\sqrt{1+q^2}$ and $\|X_2\|_{L^2(P_2)}=\sqrt Q$. Integration gives $\int r_G\le C B_s+Cq^2(t_b-t_0)$.
\end{proof}

\section{Correct partition and remaining initialized obligations}
One can use the following fixed groups without charging the same output twice. The initial right half-circle has a safe strong core and a mass-charged complement. The whole initial Q2 sector remains in the explicit-output ledger. Within it, the fixed bulk subcohort $\Clk$ controls $I$; the recruited strong-chart secondary subset is handled by Corollary \ref{inc:AN03:bulkclock:v1:secondary}; the remaining primary transit labels, including $\mathcal M_q$, need their own energy comparison. One may either include those transit labels in a separately named primary bulk or carry an additional $O(q^{1/4}H)$ term if that transport bound is proved. Initial Q3/background labels retain their pre-existing separate obligations.

The sharp quadratic secondary exponent is $\gamma$, and its small fixed-power-loss version $\Gamma_\nu>3/10$ is compatible with the already reviewed right-half-circle margin $\Delta_\nu>3/25$ at $\kappa=10$, $\nu\le1/100$. These are conditional compatible powers, not initialized applicability. Charging a possible secondary mass $O(q^{\Gamma_\nu}H)$ by the generic $1/q$ bound would give an unusable upper budget; the componentwise energy calculation avoids that particular loss. It does not prove that the actual secondary field diverges under another decomposition.

The common missing radial estimate is still the signed original-flow capped-gain inequality, including the strong diagonal contribution, the c-form all-label distance envelope through exit, and the radial defect along escaping labels. This increment proves how that package controls secondary weak energy. It does not supply the package. In addition, a genuine original-flow argument must transfer the joint entry profile and bulk seed, verify the residual and Cartesian entry inputs to Theorem \ref{inc:AN03:bulkclock:v1:bulkclockclosure} (or otherwise bound the bulk defect), control primary transit energy and the residual norms in Proposition \ref{inc:AN03:bulkclock:v1:residualinputs}, justify any amplitude cap using the exact response/antisymmetric/negative-gate ledger, and retain selected-residual and background errors in the field reduction. Final $q^2$-scale probe transfer is unchanged as a separate requirement.

\begin{thebibliography}{9}\small
\bibitem{CLOCK}\path{AN03_exact_amplitude_clock_and_c_passage_v1.tex}. Labels \texttt{identity}, \texttt{endpoint}, \texttt{crossenergy}, \texttt{passage}, \texttt{forcing}, prefix \texttt{inc:AN03:clock:v1:}.
\bibitem{Q2E}\path{AN02_Q2_boundary_layers_and_cohort_energy_v1.tex}. Labels \texttt{crossing}, \texttt{uncrossed}, \texttt{profile}, \texttt{energy}, and equation \texttt{axisU1upper}, prefix \texttt{inc:AN02:q2completion:v1:}.
\bibitem{TAIL}\path{AN03_mass_charged_tail_and_selection_defects_v1.tex}. Capped-gain integration and rational margin, prefix \texttt{inc:AN03:tailrepair:v1:}.
\bibitem{PROFILE}\path{AN02_strong_entry_profile_and_tail_budget_v1.tex}. Exact-center profile and original-state outside mass estimate, prefix \texttt{inc:AN02:profile:v1:}.
\bibitem{WEAK}\path{AN03_weak_transit_forcing_and_orbit_selection_v1.tex}. Leading-system weak return and selected-orbit comparison; no original-flow capture premise is imported.
\end{thebibliography}
\end{document}
